\section{Background}
\label{sec:background}

% Written 2026-09-17 at the user's request. Scope chosen to cover exactly what the rest of the paper
% assumes and no more: what a semantics-preserving transform is, what the two prediction directions
% are, and the three adaptation mechanisms the results section compares. Every citation here was
% verified against arxiv.org rather than recalled; Granite has no arXiv paper for its 3.x instruct
% family and is cited by release.

\subsection{Semantics-preserving obfuscation}
An obfuscating transform rewrites a program's surface while preserving its observable behaviour:
for every input, the transformed program returns what the original returned. The standard taxonomy
separates \emph{lexical} transforms, which rewrite identifiers and strip the names and annotations a
reader relies on, from \emph{control} transforms, which restructure the flow the computation takes,
and \emph{data} transforms, which change how values are represented and
computed~\cite{collberg1997taxonomy,collberg2002manufacturing}. The distinction matters here because
it predicts what a model must recover. A lexical transform destroys evidence the model may have been
relying on but leaves the computation's shape intact; a control transform leaves the identifiers and
changes the shape; a data transform changes the arithmetic itself. An adaptation that generalises
across all three has learned something about the computation. One that generalises within a family
has learned that family's surface.

Two properties of the class are what make it a usable instrument. First, the transforms
\emph{compose}: an obfuscator may rename identifiers inside flattened control flow inside rewritten
expressions, and real ones do. Composition is not a harder version of the same test but a different
one, since a model may handle each part alone and still fail the stack. Second, because the
transforms preserve semantics by construction, the ground truth is computable from the clean parent
and is independent of any transform --- so a grader can be exact, and a claim of invariance is
falsifiable rather than a matter of judgement.

\subsection{Predicting behaviour in two directions}
We evaluate program understanding as behaviour prediction rather than as generation or recovery.
Given a program and an input, \emph{output prediction} asks what it returns; given a program and a
return value, \emph{input prediction} asks for a call that produces it. Both are standard, and
current code models are markedly imperfect at both even on unobfuscated
code~\cite{cruxeval}. The pairing is what we use: the two directions share a program, a gold value
and a grader, so an adaptation that improves one and damages the other has not become better at
reading the program. Input prediction also has a property the forward direction lacks --- inversion
is many-to-one, so grading it requires executing the produced call rather than matching it, which we
do throughout.

We never train on, or evaluate, recovery of the original source. That is the clean separation from
the deobfuscation lineage: a model that can undo a transform tells us about the transform, whereas a
model that answers correctly on the transformed program tells us about its understanding of the
computation.

\subsection{Adapting a model, and composing adaptations}
Fine-tuning a code model on a new distribution is now routinely done with low-rank
adaptation~\cite{lora}, which freezes the base weights and trains a pair of low-rank matrices per
target module. The practical consequence for this paper is not efficiency but \emph{modularity}:
each adapter is a small, separable object, so one can be trained per transform and the resulting set
composed afterwards --- which is the design space the results section explores.

Three ways to compose them appear in the literature and are the baselines we compare.
\textbf{Pooled training} ignores the modularity and trains a single adapter on the union of all the
data; it is the obvious thing to do and the one most likely to be done in practice.
\textbf{Routing} keeps the adapters separate and selects or mixes among them per input, which needs a
gate at inference and pays for every expert it keeps loaded. \textbf{Weight-space merging} collapses
the adapters into one set of weights before inference, so it costs nothing at run time. The merging
methods we use are the two standard ones. TIES resolves the interference between task vectors by
trimming small updates and, where adapters disagree on a parameter's sign, keeping only those that
agree with the majority~\cite{ties}. DARE adds a preliminary step that randomly drops a fraction of
the delta entries and rescales the survivors so the expected update is preserved, on the finding that
fine-tuning deltas are highly redundant~\cite{dare}. Both are used unmodified; the contribution of
this paper is not a merging algorithm but an evaluation that separates these three mechanisms where
they actually differ.
